\documentclass[11pt]{article}

% ---------- Packages ----------
\usepackage[a4paper,margin=2.2cm]{geometry}
\usepackage{graphicx}
\usepackage{booktabs}
\usepackage{longtable}
\usepackage{array}
\usepackage{xcolor}
\usepackage{titlesec}
\usepackage{enumitem}
\usepackage{hyperref}
\usepackage{fancyhdr}
\usepackage{listings}
\usepackage{caption}
\usepackage{parskip}
\usepackage{amsmath}

% ---------- Colors ----------
\definecolor{headingblue}{HTML}{2C4F6B}
\definecolor{lightgrey}{HTML}{F1EFE8}
\definecolor{codebg}{HTML}{F7F7F5}

% ---------- Section styling ----------
\titleformat{\section}{\large\bfseries\color{headingblue}}{\thesection}{0.6em}{}
\titleformat{\subsection}{\normalsize\bfseries\color{headingblue}}{\thesubsection}{0.6em}{}
\titleformat{\subsubsection}{\normalsize\bfseries}{\thesubsubsection}{0.6em}{}
\titlespacing*{\section}{0pt}{1.6ex plus 0.4ex minus .2ex}{1ex plus .2ex}
\titlespacing*{\subsection}{0pt}{1.2ex plus 0.3ex minus .2ex}{0.8ex plus .2ex}

% ---------- Header/footer ----------
\pagestyle{fancy}
\fancyhf{}
\fancyhead[L]{\small Customer Risk Segmentation --- Cluster Analysis}
\fancyhead[R]{\small\thepage}
\renewcommand{\headrulewidth}{0.4pt}

% ---------- Hyperlinks ----------
\hypersetup{
  colorlinks=true,
  linkcolor=headingblue,
  urlcolor=headingblue,
  citecolor=headingblue,
  pdftitle={Customer Risk Segmentation via Cluster Analysis},
  pdfauthor={Clare Tenewaa Duku}
}

% ---------- Code listing style (R) ----------
\lstdefinestyle{rcode}{
  language=R,
  backgroundcolor=\color{codebg},
  basicstyle=\ttfamily\scriptsize,
  commentstyle=\color{gray},
  keywordstyle=\color{headingblue}\bfseries,
  stringstyle=\color{teal!70!black},
  numbers=left,
  numberstyle=\tiny\color{gray},
  numbersep=6pt,
  breaklines=true,
  breakatwhitespace=false,
  frame=single,
  rulecolor=\color{gray!40},
  framexleftmargin=6pt,
  xleftmargin=14pt,
  showstringspaces=false,
  tabsize=2
}

% ---------- Table row color helper ----------
\newcolumntype{L}[1]{>{\raggedright\arraybackslash}p{#1}}
\newcolumntype{C}[1]{>{\centering\arraybackslash}p{#1}}

\begin{document}

% ======================================================================
% COVER PAGE
% ======================================================================
\begin{titlepage}
\centering
\vspace*{3.5cm}

{\huge\bfseries Customer Risk Segmentation via Cluster Analysis \par}
\vspace{0.5cm}
{\large\itshape\color{gray} K-means and DBSCAN on a Simulated Retail Banking Dataset \par}

\vspace{2.5cm}
{\Large\bfseries Clare Tenewaa Duku \par}
\vspace{0.3cm}
{\normalsize\color{gray} \href{https://github.com/claretenewaa/Cluster-Analysis-for-Customer-Risk-Segmentations-}{github.com/claretenewaa/Cluster-Analysis-for-Customer-Risk-Segmentations-} \par}

\vspace{2cm}
{\normalsize\color{gray} Dataset: 100,300 synthetic customer records (100,000 unique customers), fixed random seed 602 for reproducibility \par}

\vfill
{\small\color{gray} Personal data science portfolio project \par}
\end{titlepage}

% ======================================================================
% 1. INTRODUCTION
% ======================================================================
\section{Introduction and Segmentation Question}

Modern retail banks routinely analyse customer-level financial data to support strategic decisions, including marketing, retention, and risk management. This report applies unsupervised cluster analysis to a simulated retail banking dataset of 100{,}000 customers to address the following segmentation question:

\begin{quote}
\bfseries ``Can customers be grouped into distinct, actionable risk profiles based on their income, savings, debt, and repayment behaviour, in order to support differentiated credit risk management and collections strategy?''
\end{quote}

The dataset was generated using a fixed random seed (602) via a supplied R simulation function, producing demographic, product-holding, transactional, credit-risk, and engagement variables, together with realistic data-quality issues (missing values, entry errors, and duplicate records) that were identified and corrected as part of this analysis (Section~\ref{sec:cleaning}).

Two complementary clustering algorithms were applied: K-means, to establish a small number of broad, business-interpretable risk segments; and DBSCAN, a density-based method used to identify anomalous, high-risk customers that a centroid-based method would otherwise absorb into an average cluster. The remainder of this report documents variable selection, data preparation, methodology, results, cluster interpretation, business recommendations, a comparison of the two methods, and limitations.

% ======================================================================
% 2. VARIABLE SELECTION
% ======================================================================
\section{Variable Selection and Justification}

\subsection{Selected Variables}

Nine variables were selected for the cluster analysis, chosen to jointly represent three complementary dimensions of financial risk:

\begin{itemize}[leftmargin=1.4em]
  \item \textbf{Repayment capacity}: \texttt{monthly\_income}, \texttt{savings\_balance}, \texttt{current\_balance} --- buffers against financial shocks.
  \item \textbf{Leverage}: \texttt{total\_outstanding\_debt}, \texttt{debt\_to\_income\_ratio}, \texttt{num\_active\_loans} --- current indebtedness relative to means.
  \item \textbf{Repayment behaviour}: \texttt{num\_missed\_payments}, \texttt{overdraft\_usage\_yearly}, \texttt{credit\_score} --- realised repayment conduct.
\end{itemize}

Combining capacity, leverage, and behaviour avoids a one-dimensional view of risk and allows segments to separate, for example, ``high income but high leverage'' customers from ``low income but low leverage, reliable payers'' --- a distinction a single variable could not draw.

\subsection{Excluded Variables}

Demographic fields (age, gender, region, education, marital status) were excluded as they are not direct financial-risk indicators and their use in risk segmentation raises proxy-discrimination concerns. Identifier fields (\texttt{customer\_id}, \texttt{branch\_id}, \texttt{relationship\_manager\_id}) were dropped as meaningless for clustering. Engagement and service variables (\texttt{satisfaction\_score}, \texttt{branch\_visits\_yearly}, \texttt{num\_complaints}, \texttt{churned}, \texttt{has\_mobile\_banking}) were excluded as they reflect service experience rather than credit risk. \texttt{loan\_default\_flag} was excluded because it is the eventual supervised-learning target; including it would leak the outcome into an unsupervised risk-clustering exercise.

\subsection{Variables Rejected for Correlation / Redundancy}

\texttt{has\_personal\_loan} and \texttt{has\_mortgage} were dropped because they are already summarised in \texttt{num\_active\_loans}. \texttt{avg\_transaction\_value} and \texttt{avg\_monthly\_transactions} were excluded as they correlate strongly with \texttt{monthly\_income} and mainly reflect channel activity rather than risk.

% ======================================================================
% 3. DATA CLEANING & PREPROCESSING
% ======================================================================
\section{Data Cleaning and Preprocessing}
\label{sec:cleaning}

\subsection{Identified Data Quality Issues}

Three distinct issues were identified by direct inspection of the nine selected variables in the generated dataset.

\subsubsection*{Issue 1 --- Exact duplicate records}
300 fully duplicated rows (0.3\% of 100{,}300 records) were found, consistent with a data-integration duplication error. \textbf{Remedy:} removed via \texttt{distinct()} row de-duplication. This was preferred over fuzzy matching because the rows were byte-for-byte identical, so exact removal loses no information and introduces no bias.

\subsubsection*{Issue 2 --- Implausible / impossible values (entry errors)}
\begin{itemize}[leftmargin=1.4em]
  \item \texttt{monthly\_income}: 176 values were non-positive (including $-500$) or implausibly extreme (up to 5{,}000{,}000, against a clean 99.9th percentile of $\sim$29{,}068). Recoded to \texttt{NA} and handled with Issue 3.
  \item \texttt{debt\_to\_income\_ratio}: 25 negative values (a ratio cannot be negative). Corrected via \texttt{abs()}, which exactly reverses a sign-flip entry error.
  \item \texttt{savings\_balance}: $\sim$500 values inflated roughly 1000$\times$, consistent with an extra-zero entry error. Winsorised (capped) at the 99.5th percentile rather than deleted, preserving the customer record while preventing a handful of rows from dominating distance-based clustering.
\end{itemize}
Capping/correction was preferred over deletion in each case because the affected customers likely have genuine, otherwise-usable records; deleting them would needlessly shrink the sample and could bias segments if the errors are not random.

\subsubsection*{Issue 3 --- Missing values}
\texttt{monthly\_income} was missing in 6.0\% of records and \texttt{credit\_score} in 4.0\%. \textbf{Remedy:} median imputation, with a companion \texttt{*\_was\_missing} flag retained for each variable. Median imputation was preferred over mean imputation because both variables are heavily right-skewed, so the mean is pulled upward by extreme earners while the median better represents a typical customer. Listwise deletion was rejected as it would discard roughly 10\% of customers combined and could bias segments if missingness is related to customer type (e.g.\ unverified self-employed income).

\subsection{Preprocessing: Transformation and Scaling}

Skewness and spread were checked empirically for each variable before choosing a scaling method, rather than applying one method uniformly.

\begin{table}[h]
\centering
\small
\begin{tabular}{L{3.1cm} L{4.3cm} L{2.6cm} L{3.2cm}}
\toprule
\textbf{Variable group} & \textbf{Variables} & \textbf{Skew (raw $\to$ log)} & \textbf{Method chosen} \\
\midrule
Monetary (right-skewed) & income, savings, current balance, debt, DTI & up to 25.6 $\to$ $-0.6$ to 4.2 & log1p, then Z-score \\
Bounded, symmetric spread & credit\_score & low & Robust scaling (median/IQR) \\
Zero-inflated counts & active loans, missed payments, overdraft use & median = IQR = 0 & Min-Max normalisation \\
\bottomrule
\end{tabular}
\caption{Preprocessing method chosen per variable group.}
\end{table}

The five monetary variables were log1p-transformed to compress heavy right tails (a handful of wealthy customers would otherwise dominate Euclidean-distance clustering), then Z-scored so all five carry comparable variance. \texttt{credit\_score} has a genuinely symmetric spread with non-zero IQR, so median/IQR robust scaling was used directly. The three count variables (\texttt{num\_active\_loans}, \texttt{num\_missed\_payments}, \texttt{overdraft\_usage\_yearly}) are heavily zero-inflated --- checking directly confirmed their median \emph{and} IQR are both 0 ($\geq$75\% of customers have zero missed payments or overdraft use), which makes robust scaling mathematically degenerate; Min-Max normalisation was used instead, as it bounds each variable to $[0,1]$ without relying on a spread statistic. \texttt{debt\_to\_income\_ratio} and \texttt{total\_outstanding\_debt} retain moderate residual skew even after logging, reflecting a structural point mass at zero (customers with no debt) rather than a preprocessing failure --- noted as a limitation in Section~\ref{sec:limitations}.

% ======================================================================
% 4. METHODOLOGY
% ======================================================================
\section{Methodology}

\subsection{Clustering Algorithms Employed}

Two complementary algorithms were applied to the preprocessed 9-feature, 100{,}000-row matrix:

\begin{itemize}[leftmargin=1.4em]
  \item \textbf{K-means} --- a centroid-based partitioning method, used to define a small number of broad, interpretable customer segments for the risk-classification business objective.
  \item \textbf{DBSCAN} --- a density-based method, used as the second algorithm to detect anomalous, high-risk customers as explicit ``noise'' rather than forcing them into an averaged cluster.
\end{itemize}

\subsection{Choice of Parameters}

K-means was run with \texttt{nstart = 10} (10 random centroid initialisations per $k$, keeping the best solution) and a maximum of 50 iterations, for $k = 1$ to $10$, using \texttt{set.seed(602)} for reproducibility.

For DBSCAN, \texttt{minPts} was set to 18 (twice the feature dimensionality, a standard rule of thumb for continuous, moderate-dimensional data). \texttt{eps = 0.60} was chosen from the $k$-distance elbow: the sorted 18th-nearest-neighbour distance curve rises gently up to roughly the 90th percentile (0.568) before accelerating sharply beyond the 96th percentile (0.725+); 0.60 sits just past the knee, balancing a meaningful density threshold against an excessive noise fraction.

\subsection{Elbow and Silhouette Analyses}

\begin{figure}[h]
\centering
\includegraphics[width=0.72\linewidth]{figures/final_elbow_plot.png}
\caption{K-means elbow method --- total within-cluster sum of squares (WSS) by $k$.}
\end{figure}

Total WSS declines steeply from $k=2$ to $k=5$ (19--22\% per additional cluster) and flattens markedly from $k=6$ onward (7--9\% per step) --- the elbow sits at $k=5$ or $6$.

\begin{figure}[h]
\centering
\includegraphics[width=0.72\linewidth]{figures/final_silhouette_plot.png}
\caption{K-means silhouette method --- average silhouette width by $k$ (evaluated on an 8{,}000-customer sample; silhouette computation is $O(n^2)$ and infeasible on the full 100{,}000 rows).}
\end{figure}

The global maximum average silhouette width occurs at $k=2$ (0.293), but a clear local peak appears at $k=6$--$7$ (0.286--0.288), nearly as high, following a dip at $k=3$.

\subsubsection*{Reconciling the two diagnostics}
The two methods nominally disagree (elbow $\approx 5$--6; silhouette's global optimum $=2$), so both were checked against the business objective before a final choice was made. $k=2$ is statistically the tightest split, but profiling shows it merely separates ``low outstanding debt'' (71\% of customers) from ``high debt / elevated DTI'' (29\%) --- too coarse to support differentiated risk strategy. $k=6$ sits at a local silhouette peak nearly matching the global maximum, coincides with where the elbow's marginal gains flatten, and --- critically --- produces six segments that are individually interpretable and include a distinct, small, high-risk cluster (Section~\ref{sec:profiles}). $k=6$ was therefore selected as the final solution, prioritising business actionability where the two purely statistical diagnostics were closely matched.

% ======================================================================
% 5. RESULTS
% ======================================================================
\section{Results}

\subsection{Cluster Visualisations}

No dendrogram is included, as hierarchical clustering was not one of the two algorithms used for this analysis (K-means and DBSCAN were chosen instead, for the reasons discussed in Section~4.1).

\begin{figure}[h]
\centering
\includegraphics[width=0.68\linewidth]{figures/cluster_sizes.png}
\caption{K-means ($k=6$) cluster sizes.}
\end{figure}

\begin{figure}[h]
\centering
\includegraphics[width=0.72\linewidth]{figures/final_kmeans_pca.png}
\caption{K-means ($k=6$) clusters projected onto the first two principal components (PC1 = 34.1\%, PC2 = 28.8\% of variance; 12{,}000-customer plotting sample).}
\end{figure}

\begin{figure}[h]
\centering
\includegraphics[width=0.72\linewidth]{figures/final_dbscan_pca.png}
\caption{DBSCAN clusters projected onto the same PCA axes; grey points are density-based outliers (``noise'').}
\end{figure}

\clearpage
\subsection{Summary Tables}

\subsubsection*{K-means ($k=6$) cluster profile (raw-scale means)}
\begin{table}[h]
\centering
\small
\begin{tabular}{c c c c c c c c c}
\toprule
\textbf{Cl.} & \textbf{n (\%)} & \textbf{Income} & \textbf{Savings} & \textbf{Debt} & \textbf{DTI} & \textbf{Credit} & \textbf{Missed} & \textbf{OD/yr} \\
\midrule
1 & 14350 (14.3\%) & \$6942.83  & \$56877.49 & \$38713.60  & 0.47  & 701.51 & 0.06 & 0.56 \\
2 & 26947 (26.9\%) & \$5458.54  & \$36607.79 & \$2.85      & 0.00  & 707.18 & 0.01 & 0.01 \\
3 & 4059 (4.1\%)   & \$1980.51  & \$11390.35 & \$147483.04 & 11.57 & 661.09 & 2.27 & 0.21 \\
4 & 20444 (20.4\%) & \$1108.34  & \$3119.28  & \$3.27      & 0.00  & 647.88 & 0.13 & 0.01 \\
5 & 10510 (10.5\%) & \$2080.83  & \$6853.64  & \$8984.40   & 0.48  & 641.95 & 0.24 & 0.94 \\
6 & 23690 (23.7\%) & \$5080.07  & \$40889.19 & \$3.60      & 0.00  & 702.21 & 0.02 & 0.01 \\
\bottomrule
\end{tabular}
\caption{K-means cluster profile on raw (unscaled) variables. OD/yr = overdraft usage per year.}
\end{table}

\subsubsection*{DBSCAN: noise (outlier) vs.\ clustered customers (raw-scale means)}
\begin{table}[h]
\centering
\small
\begin{tabular}{L{4.2cm} C{3.2cm} C{3.2cm}}
\toprule
\textbf{Variable} & \textbf{Noise (n=3,303)} & \textbf{Clustered (n=96,697)} \\
\midrule
monthly\_income          & 2766.63 & 4245.19 \\
savings\_balance         & 43539.90 & 29055.20 \\
current\_balance         & 2407.31 & 2359.76 \\
total\_outstanding\_debt & 80499.26 & 10165.15 \\
debt\_to\_income\_ratio  & 9.28 & 0.29 \\
credit\_score            & 633.06 & 686.09 \\
num\_active\_loans       & 1.08 & 0.34 \\
num\_missed\_payments    & 2.03 & 0.10 \\
overdraft\_usage\_yearly & 0.90 & 0.17 \\
\bottomrule
\end{tabular}
\caption{DBSCAN noise points vs. clustered customers on raw risk indicators.}
\end{table}

% ======================================================================
% 6. CLUSTER PROFILES & BUSINESS LABELS
% ======================================================================
\clearpage
\section{Cluster Profiles and Business Labels}
\label{sec:profiles}

\subsubsection*{Cluster 2 --- ``High-Income Savers, No Debt'' (26.9\%, n=26,947)}
Highest credit score (707), zero debt, no current account, income \$5{,}459. The largest, lowest-risk segment: financially comfortable customers who save rather than borrow. \textbf{Business label: Minimal risk}

\subsubsection*{Cluster 6 --- ``Affluent Transactors, No Debt'' (23.7\%, n=23,690)}
Similar to Cluster 2 but distinguished by active current-account usage (highest current balance, \$6{,}594) alongside zero debt and strong credit (702). Engaged, debt-free, day-to-day banking relationship. \textbf{Business label: Minimal risk}

\subsubsection*{Cluster 4 --- ``Underbanked, Low-Income'' (20.4\%, n=20,444)}
Lowest income (\$1{,}108) and lowest savings (\$3{,}119), essentially no debt, but the weakest credit score (648) of any segment. Low risk today mainly because these customers carry no debt, not because of strong credit fundamentals. \textbf{Business label: Low risk, but fragile}

\subsubsection*{Cluster 1 --- ``Affluent Leveraged (Mortgage-Holders)'' (14.3\%, n=14,350)}
High income (\$6{,}943) and savings (\$56{,}877) alongside meaningful debt (\$38{,}714, DTI 0.47) --- consistent with mortgage or large personal-loan holders. Good credit (702), very low missed payments (0.06). \textbf{Business label: Low-moderate risk}

\subsubsection*{Cluster 5 --- ``Overdraft-Reliant, Moderate Risk'' (10.5\%, n=10,510)}
Modest income (\$2{,}081), one of the lowest credit scores (642), and by far the highest overdraft usage (0.94/year). Debt is moderate (DTI 0.48) but repeated overdraft reliance signals cash-flow stress. \textbf{Business label: Moderate risk}

\subsubsection*{Cluster 3 --- ``Financially Distressed'' (4.1\%, n=4,059)}
The clear high-risk segment: low income (\$1{,}981) combined with very high debt (\$147{,}483) and an extreme debt-to-income ratio (11.57 --- roughly 25$\times$ the next-highest cluster). Highest missed-payment rate (2.27/year) of any segment. \textbf{Business label: High risk}

% ======================================================================
% 7. INTERPRETATION, RECOMMENDATIONS, COMPARISON
% ======================================================================
\section{Interpretation, Recommendations, and Comparison of Methods}

\subsection{Interpretation and Business Recommendations}

\begin{itemize}[leftmargin=1.4em]
  \item Cluster 3 (Financially Distressed, 4.1\%) should be prioritised for proactive risk management: early-warning monitoring, restructuring or hardship-programme outreach, and tighter approval thresholds on further lending given the extreme DTI and missed-payment rate.
  \item Cluster 5 (Overdraft-Reliant, 10.5\%) is a moderate-risk, high-value intervention target: targeted cash-flow products (e.g.\ graduated overdraft limits, budgeting tools, or a consolidation loan) could reduce reliance on high-cost overdraft facilities before customers migrate toward Cluster 3.
  \item Clusters 2 and 6 (jointly 50.6\%, debt-free, high credit) are prime candidates for cross-selling accretive products --- credit cards, investment accounts, fixed deposits --- since capacity is high and current risk exposure is minimal.
  \item Cluster 1 (Affluent Leveraged, 14.3\%) represents the bank's mortgage/large-loan relationship base; retention and refinancing offers are more relevant here than risk mitigation, given strong credit and low missed payments despite high debt.
  \item Cluster 4 (Underbanked, 20.4\%) is low-risk today but financially fragile (weak credit score, minimal buffers); appropriate for financial-inclusion products (starter credit-builder cards, low-fee savings) rather than collections attention.
\end{itemize}

\subsection{Comparison of Clustering Methods}

K-means and DBSCAN answered different but complementary parts of the segmentation question. K-means, by construction, assigns every customer to one of a fixed number of roughly spherical clusters, which produced six broad, evenly-sized, business-interpretable segments --- well suited to designing differentiated marketing and relationship-management strategies for the bulk of the customer base.

DBSCAN made no assumption about cluster shape or count, and instead of forcing every customer into the nearest centroid, it explicitly separated 3.3\% of customers (3{,}303 records) as density-based ``noise''. Comparing this noise set against the rest of the population on every raw risk variable (Table, Section 5.2) shows it is dramatically worse on every dimension simultaneously --- 32$\times$ higher DTI, 8$\times$ more outstanding debt, and 20$\times$ more missed payments than the clustered majority. This is a materially higher concentration of risk than even K-means' own worst cluster (Cluster 3), because DBSCAN's noise definition specifically isolates points that are sparse in the full 9-dimensional feature space rather than merely distant from a single centroid.

In practice, the two methods are best used together: K-means for the primary segmentation used in day-to-day marketing and relationship strategy, and DBSCAN as a targeted overlay to flag the small number of customers who most urgently warrant manual underwriting review or early intervention --- a distinction K-means alone cannot make, since its centroid-based objective necessarily averages outliers into their nearest cluster rather than isolating them.

% ======================================================================
% 8. LIMITATIONS & CONCLUSION
% ======================================================================
\section{Limitations and Conclusion}
\label{sec:limitations}

\subsection{Limitations}

\begin{itemize}[leftmargin=1.4em]
  \item \textbf{Simulated data}: relationships between variables (e.g.\ income $\to$ debt $\to$ default probability) were generated by parametric formulas rather than observed from real customer behaviour; cluster structure may not fully generalise to a live portfolio.
  \item \textbf{Residual skew}: \texttt{debt\_to\_income\_ratio} and \texttt{total\_outstanding\_debt} retained moderate skew even after log-transformation, due to a structural point mass at zero (no-debt customers) that a single continuous transform cannot fully resolve; a two-part (zero-inflated) model was outside this project's scope.
  \item \textbf{Silhouette evaluated on a sample}: due to $O(n^2)$ computational cost, silhouette width was computed on an 8{,}000-customer sample rather than the full 100{,}000 rows; while cluster assignments themselves came from the full-data fit, the reported silhouette values carry sampling variability.
  \item \textbf{DBSCAN parameter sensitivity}: \texttt{eps} and \texttt{minPts} were chosen via a $k$-distance elbow heuristic and a standard rule of thumb, respectively; different choices materially affect the noise fraction and cluster count, and were not exhaustively grid-searched.
  \item \textbf{Static, cross-sectional segmentation}: clusters reflect a single snapshot of each customer's financial position and do not capture trajectory (e.g.\ a customer moving from Cluster 2 toward Cluster 5 over time), which would strengthen early-warning value.
\end{itemize}

\subsection{Conclusion}

This analysis identified six actionable K-means risk/value segments and a DBSCAN-based outlier overlay isolating the most acute risk cases, using nine variables spanning income, leverage, and repayment behaviour selected specifically for the risk-classification business objective. The two methods were shown to be complementary rather than redundant: K-means for portfolio-wide segmentation strategy, and DBSCAN for surfacing the small minority of customers requiring urgent, individualised risk attention. Together they provide a practical foundation for differentiated marketing, retention, and collections strategy across the simulated customer base.

% ======================================================================
% 9. APPENDIX
% ======================================================================
\clearpage
\appendix
\section{Appendix --- Fully Commented R Code}

The complete, executable R pipeline (data generation, variable selection, cleaning, preprocessing, and cluster analysis) is reproduced below, extracted directly from the project's R Markdown source. Comment lines beginning with \texttt{\#'} correspond to the narrative sections of the analysis; standard \texttt{\#} comments mark individual processing steps.

\lstinputlisting[style=rcode]{appendix_code.R}

\end{document}
